\section{The rule in practice, and who splits}
\label{sec:practice}

\begin{table*}[t]
\centering
\small
\begin{adjustbox}{max width=\linewidth}\begin{tabular}{llcccccc}
\toprule
Data & Model & Recall & Precision & F1 & Tokens & Recall & Cost \\
 & & whole $\to$ rule & whole $\to$ rule & whole $\to$ rule & rule/whole & rule $-$ 44 px & $/C_{\min}$ \\
\midrule
LegendSpot & GPT-6.1 & 0.97 $\to$ 1.00 & 0.93 $\to$ 0.97 & 0.95 $\to$ 0.98 & 8.6 & +0.00 & 3.5 \\
LegendSpot & Claude Opus 5.5 & 0.83 $\to$ 0.92 & 0.82 $\to$ 0.92$^{*}$ & 0.81 $\to$ 0.90 & 5.2 & $-$0.05 & 2.6 \\
LegendSpot & Gemini 3.8 & 0.62 $\to$ 0.90$^{*}$ & 0.72 $\to$ 0.88$^{*}$ & 0.65 $\to$ 0.87$^{*}$ & 7.6 & $-$0.04 & 1.6 \\
LegendSpot & Gemini 4 & 0.85 $\to$ 0.98$^{*}$ & 0.89 $\to$ 0.93 & 0.85 $\to$ 0.94$^{*}$ & 7.5 & +0.00 & 1.6 \\
LegendSpot & Qwen3-VL & 0.37 $\to$ 0.56$^{*}$ & 0.75 $\to$ 0.93$^{*}$ & 0.37 $\to$ 0.56$^{*}$ & 6.1 & $-$0.07 & 2.0 \\
\midrule
Floor plans & GPT-6.1 & 0.63 $\to$ 0.61 & 0.80 $\to$ 0.84 & 0.63 $\to$ 0.64 & 4.4 & n/a & 3.3 \\
Floor plans & Claude Opus 5.5 & 0.51 $\to$ 0.56 & 0.81 $\to$ 0.79 & 0.57 $\to$ 0.59 & 3.1 & n/a & 2.0 \\
Floor plans & Gemini 3.8 & 0.25 $\to$ 0.43$^{*}$ & 0.58 $\to$ 0.78$^{*}$ & 0.32 $\to$ 0.50$^{*}$ & 6.8 & $-$0.13$^{*}$ & 1.6 \\
Floor plans & Gemini 4 & 0.50 $\to$ 0.57 & 0.70 $\to$ 0.83$^{*}$ & 0.56 $\to$ 0.63$^{*}$ & 6.7 & $-$0.01 & 1.6 \\
Floor plans & Qwen3-VL & 0.05 $\to$ 0.12$^{*}$ & 0.28 $\to$ 0.46$^{*}$ & 0.06 $\to$ 0.15$^{*}$ & 3.5 & $-$0.11$^{*}$ & 2.0 \\
\bottomrule
\end{tabular}
\end{adjustbox}
\caption{The rule ($S^\ast = 1.5$, without step 3 of \cref{alg:rule}) against the whole-image call on 34 LegendSpot drawings and 60 floor-plan mosaics. Recall, precision and F1: means per image under one-to-one matching; duplicates in overlaps are not merged and count against precision, which omits images where either call returned no point (Qwen3-VL: 17 and 47; others at most 5). Tokens: billed input tokens (prompts and reference image included), rule over whole-image call. Recall $-$ 44 px: difference to the 44 px recipe (n/a: not run). Cost $/\Cmin$: the rule's image tokens over $\Cmin$, median. ${}^\ast$: 95\% paired bootstrap interval excludes 0. Gemini 4's whole-image call asks for 0 to 1000 coordinates (\supp{B}).}
\label{tab:main}
\end{table*}

The rule raises recall most where the whole image gives targets the fewest tokens, at a cost known in advance, and zoom agents that succeed do so by splitting the image themselves. We compare it with the whole-image call and the 44 px recipe for GPT-6.1, Claude Opus 5.5, Gemini 3.8, Gemini 4 and Qwen3-VL, \CtCells{} model and data settings in all.

\paragraph{Recall and cost (\cref{tab:main}).} Recall rose significantly in \CtSigRec{} settings: Gemini 3.8 on LegendSpot from \CtLsGemThreeEightRecW{} to \CtLsGemThreeEightRecR{} (\CtLsGemThreeEightRel\% relative), Gemini 4 on LegendSpot from \CtLsGemFourRecW{} to \CtLsGemFourRecR{} (\CtLsGemFourRel\%), Qwen3-VL on LegendSpot from \CtLsQwenThreevlRecW{} to \CtLsQwenThreevlRecR{} (\CtLsQwenThreevlRel\%), Gemini 3.8 on floor-plan mosaics from \CtFpGemThreeEightRecW{} to \CtFpGemThreeEightRecR{} (\CtFpGemThreeEightRel\%), and Qwen3-VL on floor-plan mosaics from \CtFpQwenThreevlRecW{} to \CtFpQwenThreevlRecR{} (\CtFpQwenThreevlRel\%); relative gains use unrounded recalls. Precision rose significantly in \CtSigPrec{} of \CtCells{} settings, and so did F1 (\CtSigFOne{} settings). The price is \CtTokMin{} to \CtTokMax{} times the whole-image call's tokens, \CtTokDMin{} to \CtTokDMax{} times the 44 px recipe's, and \CtCminMin{} to \CtCminMax{} times $\Cmin$. Across the \CgainN{} settings, the gain falls as the whole-image call's median $S$ rises (Spearman $\rho = \CgainRhoSw$, $p = \CgainPSw$). The cost was fixed before the first call: in all \CcostPairs{} pairs of a drawing and a model, the rule sent exactly the planned views, and every view of a drawing was billed the same number of input tokens. The tokens come from raising $S$ to 1.5, and the gap to $\Cmin$ from images only a few views wide (\cref{cor:safe}); \CtOneViewGpt{} of GPT-6.1's 60 floor-plan mosaics need a single view.

\paragraph{Small gains point to step 3.} On floor plans the 44 px recipe recovers more targets than the rule for Gemini 3.8 and Qwen3-VL (\CtFpGemThreeEightdRecD{} and \CtFpQwenThreevldRecD). There the recipe cuts much smaller views, and the rule uses as little as \CtTokDMin{} of its tokens: for Qwen3-VL this mainly lowers $L$, and for the fixed-grid Gemini 3.8 a smaller view also raises $S$. On the same grid, Gemini 4, which needs much less $S$ (\cref{sec:why}), shows only \CtFpGemFourdRecD, not significant. The rule protects within an explicit factor of the lower bound; where $L$ binds, step 3 splits its views further at the same scale, at the cost given in \supp{I}.

\paragraph{Methods without a VLM, and natural images.} Template matching reaches \FTMdevAll{} F1 on the 34 drawings, and the best tiles (the 44 px recipe and its fixed-grid counterpart) score above it (Claude Sonnet 5.5 \FDsonnetFiveFiveAll, Gemini 3.8 \FDgemThreeEightAll). On the 10 further drawings whose reference symbols come from a symbol library, template matching falls to \XrdFTM{} and the VLMs are at least 0.30 higher. On FSC-147, same-scale $2 \times 2$ tiles raise small-object F1 for GPT-5.6 by \UfscODgptFiveSixsmallTU{} and for Claude Sonnet 5.5 by \UfscODsonnetFiveFivesmallTU{}: lowering $L$ at the same scale pays on natural images too.

\paragraph{Who splits: zoom agents.} At the same reasoning setting for agent and tiles, GPT-5.6's zoom agent reaches \XnFAGgptFiveSixSheets{} F1 on full sheets against \XnFDRgptFiveSixSheets{} for its tiles (\XnHOnegptFiveSix, Holm $p = \XnXNineHOnegptFiveSixHolm$); Claude Sonnet 5.5's agent reaches \XnFAGsonnetFiveFiveSheets{} against \XnFDRsonnetFiveFiveSheets, not significant. In the ablation that changes one step at a time (\cref{sec:agents}), $S$ drives GPT-5.6's gain, and native pixels add \XnHThreesonnetFiveFiveAGOVZ{} for Claude. GPT-5.6 used \XagZoomsgptFiveSix{} zooms per drawing on average and on several sheets first laid a regular grid of crops over the sheet: the agent splits the image itself, and \cref{prop:frontier} bounds its cost whenever its model's recognition is graded. Over all drawings, the Gemini 3.8 and Qwen3-VL agents fail (\XagFgemThreeEightAll{} and \XqAgF): they count as accurately as their tiles, but only \CagLocGem\% and \CagLocQwen\% of their points lie within $\tau$ of an instance, against \CtileLocGem\% and \CtileLocQwen\% for their tiles, whose points the caller maps back (\supp{J}).
